
\documentclass{article}
\usepackage{colortbl}
\usepackage{makecell}
\usepackage{multirow}
\usepackage{supertabular}

\begin{document}

\newcounter{utterance}

\centering \large Interaction Transcript for game `cladder', experiment `full\_v1.5\_default', episode 3593 with llama{-}3.1{-}8b{-}instruct.
\vspace{24pt}

{ \footnotesize  \setcounter{utterance}{1}
\setlength{\tabcolsep}{0pt}
\begin{supertabular}{c@{$\;$}|p{.15\linewidth}@{}p{.15\linewidth}p{.15\linewidth}p{.15\linewidth}p{.15\linewidth}p{.15\linewidth}}
    \# & \multicolumn{2}{c}{Player} && \multicolumn{2}{c}{Game Master} \\
    \hline

    \theutterance \stepcounter{utterance}  
    & & & \multicolumn{4}{p{0.6\linewidth}}{
        \cellcolor[rgb]{0.9,0.9,0.9}{
            \makecell[{{p{\linewidth}}}]{
                \texttt{\tiny{[P1$\langle$GM]}}
                \texttt{You are an expert in causal inference. The following question is not a typical commonsense query, but rather a meticulously designed question created by a professor specializing in causal inference, intended to assess the students' mastery of the course content.} \\
\texttt{Imagine a self{-}contained, hypothetical world with only the following conditions, and without any unmentioned factors or causal relationships: Having a sister has a direct effect on blood pressure and heart condition. Blood pressure has a direct effect on heart condition. Method 1: We look at how having a sister correlates with heart condition case by case according to blood pressure. Method 2: We look directly at how having a sister correlates with heart condition in general. To understand how having a sister affects heart condition, is it more correct to use the Method 1 than Method 2?} \\
\texttt{Start your answer with 'Yes' or 'No', followed by additional reasoning or evidence to support your explanation.} \\
            }
        }
    }
     \\ \\

    \theutterance \stepcounter{utterance}  
    & \multicolumn{4}{p{0.6\linewidth}}{
        \cellcolor[rgb]{0.561,0.737,0.561}{
            \makecell[{{p{\linewidth}}}]{
                \texttt{\tiny{[P1$\rangle$GM]}}
                \texttt{No, it is not more correct to use Method 1 than Method 2.} \\
\\ 
\texttt{In this hypothetical world, having a sister has a direct effect on blood pressure and heart condition. Blood pressure also has a direct effect on heart condition. This creates a confounding variable (blood pressure) that affects both the exposure (having a sister) and the outcome (heart condition). } \\
\\ 
\texttt{Method 1, which conditions on blood pressure, is essentially a form of stratification or matching. However, this method is still susceptible to confounding bias because it does not account for the direct effect of having a sister on blood pressure. In other words, even within the same blood pressure stratum, having a sister may still have an independent effect on heart condition.} \\
\\ 
\texttt{Method 2, on the other hand, estimates the overall effect of having a sister on heart condition without conditioning on blood pressure. This approach is more likely to capture the total effect of having a sister on heart condition, including both the direct effect and the indirect effect through blood pressure.} \\
\\ 
\texttt{In this scenario, using Method 2 is more correct because it provides a more comprehensive understanding of the causal relationship between having a sister and heart condition.} \\
            }
        }
    }
    & & \\ \\

    \theutterance \stepcounter{utterance}  
    & & & \multicolumn{2}{p{0.3\linewidth}}{
        \cellcolor[rgb]{0.9,0.9,0.9}{
            \makecell[{{p{\linewidth}}}]{
                \texttt{\tiny{[GM$|$GM]}}
                \texttt{No} \\
            }
        }
    }
    & & \\ \\

    \theutterance \stepcounter{utterance}  
    & & & \multicolumn{2}{p{0.3\linewidth}}{
        \cellcolor[rgb]{0.9,0.9,0.9}{
            \makecell[{{p{\linewidth}}}]{
                \texttt{\tiny{[GM$|$GM]}}
                \texttt{no} \\
            }
        }
    }
    & & \\ \\

    \theutterance \stepcounter{utterance}  
    & & & \multicolumn{2}{p{0.3\linewidth}}{
        \cellcolor[rgb]{0.9,0.9,0.9}{
            \makecell[{{p{\linewidth}}}]{
                \texttt{\tiny{[GM$|$GM]}}
                \texttt{game\_result = WIN} \\
            }
        }
    }
    & & \\ \\

\end{supertabular}
}

\end{document}
